% Section 4.3, from lectures/L04/README.md: the objectives, the evaluation questions, and the next
% lecture.

\section{Review}
\label{c4:sec:review}

\needspace{6\baselineskip}
\subsection*{What you should be able to do now}
After this chapter, you should:
\begin{itemize}
  \item Understand what the factory pattern is and why it is used.
  \item Understand how a factory enables switching between real drivers and stubs.
  \item Understand how the system logic can create drivers via a factory.
  \item Be able to read and follow a simple factory example.
\end{itemize}

\needspace{6\baselineskip}
\subsection*{Questions to test yourself}
\begin{enumerate}
  \item What is the purpose of the factory pattern?
  \item What does it mean that the system logic only depends on driver interfaces and not on
    concrete driver types?
  \item Who is responsible for deleting driver objects when factories return raw pointers?
  \item What is the difference between raw pointers and \code{std::unique_ptr}?
  \item Why is \code{std::unique_ptr} suitable in the factory pattern?
\end{enumerate}

\needspace{6\baselineskip}
\subsection*{Looking ahead}
Chapter~\ref{c5:ch:templates} is about templates: generic programming for embedded systems without
runtime overhead.
